\paragraph{4. Simulating robust $K$-oracle computation in $\mathsf{AM}\cap\mathsf{coAM}$.}
Throughout this section we make the following assumptions.
\begin{description}
\item[H1] $M$ is a randomized nonadaptive decision OTM.
\item[H2] $M^K$ recognizes the computable language $L$.
\item[H3] $\beta=\beta_n(q)$ is a computable closeness function.
\item[H4] $\epsilon:\mathbb N\to[0,1]$.
\item[H5] $M$ is $(\beta,\epsilon)$-robust with respect to oracle $K$.
\item[H6] $\alpha:\mathbb N\to\mathbb N$ is an unbounded non-decreasing function.
\end{description}
Also, let $k_n^M$ denote the maximum of $k_x^M$ over $x$ of length $n$.
Our overall goal is to prove the following theorem:
\paragraph{Theorem 14.}
Suppose $M,\alpha,\beta,\epsilon$ satisfy hypotheses [H1--H6]. Suppose also that $\beta_n(q)=\log(k_n^M/\epsilon(n))+2\log K(q)+\alpha(n)+5$. Then both $L$ and $\overline L$ have constant round Arthur--Merlin protocols for membership whose running time is polynomial in the running time of $M$ and $1/\epsilon$. In particular if $M$ runs in polynomial time then $L\in\mathsf{AM}\cap\mathsf{coAM}$.
The rather cumbersome condition on $\beta$ arises from the proof. It is more natural if the closeness function depends only on the query $q$ and not on the length of the input $x$ to the calling algorithm.
We prove Theorem 14 in two parts. First, in this section, we construct a parameterized family $\{J[\gamma]:\gamma\in[0,1)\}$ of computable context-sensitive oracles $J$ based on simple combinatorial properties of the machine $M$ and show that there exists an integer $\ell$ such that under [H1]--[H6], $M^{[\ell]J[\gamma]}$ computes $L$ with advantage $\epsilon/2$ for every $\gamma\in[0,1)$.
Then, in the next section, we show that if $L$ satisfies the conclusion of the first part, then it has an efficient $\mathsf{AM}$ protocol and $\mathsf{coAM}$ protocol.
\paragraph{4.1. Some additional preliminaries.}
Let $(M,x)$ be a context. As usual, for a set or quantity $Z$ that depends on the context, we denote this dependence by $Z_x^M$, but we often suppress this dependence and write simply $Z$. For example we write $r$ for $r_x^M$ and $k$ for $k_x^M$.
Let $\Gamma=\Gamma_x^M=\{(\rho,i):(\rho,i)\in\{0,1\}^r\times\{1,\ldots,k\}\}$. Under the assumption stated earlier that $k$ is a power of 2 we have that $\log|\Gamma|=r+\log k$ is an integer, and we can identify $\Gamma$ with the set $\{0,1\}^{r+\log k}$.
It will be convenient to introduce notation for some elementary functions of positive real numbers. For a nonnegative integer $s$:
\begin{itemize}
\item $\phi$ is the function given by $\phi(s)=\lceil\log(|\Gamma|/s)\rceil=\log(|\Gamma|)-\lfloor\log s\rfloor$.
\item $\mu$ is the function given by: $\mu(s)$ is the least nonnegative number $\gamma$ such that $s(1+\gamma)$ is a power of 2. Note that $\mu(s)\in[0,1)$.
\end{itemize}
It is easy to check:
\paragraph{Proposition 15.}
For $\gamma\in[0,1)$ and $s>0$, $\phi((1+\gamma)s)=\phi(s)$ if $\gamma<\mu(s)$ and $\phi((1+\gamma)s)=\phi(s)-1$ if $\gamma\geq\mu(s)$.
The following technical fact will be needed later.
\paragraph{Proposition 16.}
Let $s,t\in\mathbb N$ and $\tau\in[0,1)$ with $t\in[s/(1+\tau),s(1+\tau)]$. Let $\gamma\in[0,1)$ be such that $\gamma\notin[\mu(s)-2\tau,\mu(s)+2\tau]\cup[1-2\tau,1)\cup[0,2\tau)$. Then $\phi((1+\gamma)s))=\phi((1+\gamma)t)$.
\begin{proof}
We prove the contrapositive. Suppose $\phi((1+\gamma)s)\neq\phi((1+\gamma)t)$. Then $\lfloor\log((1+\gamma)s)\rfloor\neq\lfloor\log((1+\gamma)t)\rfloor$. This means that there is an integral power of 2 between them, which means that the interval $[(1+\gamma)s/(1+\tau),(1+\gamma)(1+\tau)s]$ contains an integral power of 2, say $2^w$. Note that $s\leq2^{w+1}<8s$. This implies that $(1+\mu(s))s=2^{w+1}$, or $(1+\mu(s))s=2^w$, or $(1+\mu(s))s=2^{w-1}$.
Case 1. Suppose $(1+\mu(s))s=2^{w+1}$. This means $s>2^w$, and since $[(1+\gamma)s/(1+\tau),(1+\gamma)(1+\tau)s]$ contains $2^w$, we have that $\tau>\gamma$. Hence $\gamma\in[0,2\tau)$.
Case 2. Suppose $(1+\mu(s))s=2^w$. Then $(1+\mu(s))s$ belongs to the interval $[(1+\gamma)s/(1+\tau),(1+\gamma)(1+\tau)s]$ so $1+\mu(s)$ belongs to the interval $[(1+\gamma)/(1+\tau),(1+\gamma)(1+\tau)]\subseteq[1+\gamma-2\tau,1+\gamma+2\tau]$ which implies $\gamma\in[\mu(s)-2\tau,\mu(s)+2\tau]$.
Case 3. Suppose $(1+\mu(s))s=2^{w-1}$. Then $2\leq2(1+\mu(s))=2^w/s\leq(1+\gamma)(1+\tau)\leq1+\gamma+2\tau$ which implies $\gamma\geq1-2\tau$.
\end{proof}
\paragraph{4.2. Replacing $K$ by the oracle $J[\gamma]$.}
As described earlier the context $(M,x)$ determines a function $q_x^M:\Gamma\to\{0,1\}^*$ where $q_x^M(\rho,i)$ is the $i$th query asked by $M$ on input $x$ when the random string is $\rho$.
For a query $q$, we define $S(q)=S_x^M(q)$ to be the set $\{(\rho,i)\in\Gamma:q_x^M(\rho,i)=q\}$.
Let $s(q)=|S(q)|$ and let $\pi(q)=\pi_x^M(q)=s_x^M(q)/|\Gamma|$. This is the probability, with respect to $(\rho,i)$ chosen uniformly from $\Gamma$ that $q(\rho,i)=q$. Note that for the function $\phi$ defined earlier we have $\phi(s(q))=\lceil\log(1/\pi(q))\rceil$.
For each $\gamma\in[0,1)$, define the function $J[\gamma]$ on queries by:
\[J[\gamma](q)=J[\gamma]_x^M(q)=\phi((1+\gamma)s(q)).\]
If $\gamma=0$ we often write $J(q)$ for $J[\gamma](q)$. From the definition we observe that for all $\gamma\in[0,1)$:
\[J[\gamma](q)\in[\log(1/\pi(q))-1,\log(1/\pi(q))+1].\tag{2}\]
We view $J[\gamma]$ as a context-sensitive oracle and we will see below that $J[\gamma]_x^M$ can be used in place of $K$ in a suitably robust oracle algorithm. The following Lemma relates $J[\gamma]$ to $K$:
\paragraph{Lemma 17.}
Let $M$ be a randomized oracle algorithm and $x$ a string. Recall that $Q_x^M$ is the set of queries $q$ such that $M$ on input $x$ asks the query $q$ with positive probability.
\begin{enumerate}
\item There are constants $C_0,C_1$ such that for any string $q\in Q_x^M$, $K(q)-J[\gamma]_x^M(q)\leq C_0K(x)+C_1$.
\item For $u>0$ let $P_x^M(u)$ be the set of strings in $Q_x^M$ such that $J[\gamma]_x^M(q)-K(q)\geq u+2\log K(q)$. Then $\pi_x^M(P_x^M(u))\leq2^{3-u}$.
\end{enumerate}
\begin{proof}
Fix $M$ and $x$. We omit the superscript $M$ and subscript $x$ from $J[\gamma],\pi,P$ and $Q$.
For the first part, order the elements of $Q$ lexicographically by the triple $(1/\pi(q),|q|,q)$. For a query $q$, let $p(q)$ be the number of $q'\in Q$ that precede $q$ in this order. Given the Turing machine $M$ and input $x$ we can construct a Turing machine of size $O(K(x))+O(1)$ that takes as input an integer $h\in\{0,\ldots,|Q|-1\}$ and outputs $q$ such that $p(q)=h$.
Therefore $K(q)\leq O(K(x))+\log p(q)+O(1)$. Since there are at most $1/\pi(q)$ members of $Q$ that precede $q$ in the order, we have $p(q)\leq(1/\pi(q))$ and therefore $\log p(q)\leq\log(1/\pi(q))\leq J[\gamma](q)+1$. Therefore $K(q)-J[\gamma](q)\leq O(K(x))+O(1)$, as required.
For the second part, (2) implies $\pi(q)\leq2^{1-J[\gamma](q)}$. For $q\in P(u)$ we therefore have $\pi(q)\leq2^{1-K(q)-u}/K(q)^2$. Summing over all $q\in P(u)$ this is at most $2^{1-u}\sum_{q\in P(u)}2^{-K(q)}/K(q)^2$ which is at most $2^{2-u}$ by Proposition 9.
\end{proof}
\paragraph{4.3. Replacing a $K$ oracle by a $J[\gamma]$ oracle.}
In this section we prove:
\paragraph{Theorem 18.}
Assume hypotheses [H1]--[H6] and that $\beta_{|x|}(q)=\log(k_x^M/\epsilon(|x|))+2\log K(q)+\alpha(|x|)+5$. Then there is an integer $\ell$ such that for all $\gamma\in[0,1]$, $M^{[\ell]J[\gamma]}(x)$ computes $L$ with advantage $\epsilon/2$, where $M^{[\ell]}$ is the OTM that outputs $L(x)$ on any input $x$ of length less than $\ell$ and otherwise executes $M$ on $x$. Letting $N=M^{[\ell]}$, there is an OTM $N$ such that $N^{J[\gamma]}(x)$ computes $L$ with advantage $\epsilon/2$.
\begin{proof}
Fix $M$. Say that $x$ is bad for $\gamma$ if $M^{J[\gamma]}(x)$ outputs $L(x)$ with probability less than $\frac12(1+\epsilon(|x|)/2)$ and let $B[\gamma]$ be the set of strings that are bad for $\gamma$. Let $B=\bigcup_{\gamma\in[0,1)}B[\gamma]$.
If $B$ is finite, let $\ell$ be the length of the largest string in $B$. Then $M^{[\ell]}$ satisfies the desired conclusion. So we assume that $B$ is infinite and derive a contradiction.
Note that $x\in B$ holds if and only if one of the infinitely many conditions $x\in B[\gamma]$ holds. The following Proposition reduces this to a finite set of conditions.
\paragraph{Proposition 19.}
Let $x\in\{0,1\}^*$ and let $\Gamma_x=\{\mu(s(q)):q\in Q_x^M\}$ (where $Q_x^M$ is the (finite) set of queries that are asked with nonzero probability by $M$ on input $x$ and $\mu$ is the function defined in the definition of $J[\gamma]$). Then $x\in B$ if and only if $x\in B[\gamma]$ for some $\gamma\in\Gamma_x$.
\begin{proof}
For each $\gamma\in[0,1)$ let $\gamma^-=\max\{\mu(s(q)):q\in Q_x^M,\mu(s(q))\leq\gamma\}$. By the definition of $J[\gamma]$, we have $J[\gamma](q)=J[\gamma^-](q)$ for all $q\in Q$. Therefore $M^{J[\gamma]}(x)$ behaves identically to $M^{J[\gamma^-]}(x)$ and so $x\in B$ if and only if there exists $\gamma\in\Gamma_x$ such that $x\in B[\gamma]$.
\end{proof}
Let us define $\gamma_x$ to be the least $\gamma\in\Gamma_x$ for which $x\in B[\gamma_x]$.
The set $\Gamma_x$ is computable from $x$, and for each $\gamma\in\Gamma_x$ we can computably determine the probability that $M^{J[\gamma]}(x)$ agrees with $L(x)$. Therefore there is an (instantiated) program of size $|M|+O(1)$ that on input $x$, tests whether $x\in B$.
Let $\beta^*(n)$ be the minimum over all $q$ and all $n'\geq n$ of $\beta_{n'}(q)$. To obtain the desired contradiction we will prove:
\paragraph{Lemma 20.}
For any string $x$ there is an oracle $K[x]$ satisfying:
\begin{itemize}
\item $K[x]$ is $\beta$-close to $K$ and therefore the probability that $M^{K[x]}$ on input $x$ differs from $M^K$ on input $x$ is at most $(1-\epsilon(|x|))/2$.
\item If $K(x)<(\beta^*(|x|)-C_1)/C_0$ (where $C_0,C_1$ are given in Lemma 17) then the probability that $M^{K[x]}$ on input $x$ differs from $M^J$ on input $x$ is at most $\epsilon(|x|)/4$.
\end{itemize}
The lemma implies any string $x$ satisfying $K(x)<(\beta^*(|x|)-C_1)/C_0$ is not in $B$, since the probability that $M^J$ on input $x$ disagrees with $M^K$ on input $x$ can be upper bounded by $(1-\epsilon(|x|)/2+\epsilon(|x|)/4\leq(1-\epsilon(|x|)/2)/2$. But by applying the second part of Proposition 10 to the language $B$ and $\kappa(n)=(\beta^*(n)-C_1)/C_0$ there is a string $x\in B$ such that $K(x)<(\beta(|x|)-C_1)/C_0$, which yields the desired contradiction. Note that $\kappa$ is an unbounded non-decreasing function as required for the application of Proposition 10 by the definition of $\beta^*$ and the fact that $\alpha$ is an unbounded non-decreasing function.
So it remains to prove Lemma 20. Given $x$, we define the oracle $K[x]$ as follows. Recall that $Q_x^M$ is the set of queries $q$ that are made by $M$ with non-zero probability when the input is $x(w)$. Let $Q_{\mathrm{close}}$ be the set of $q\in Q_x^M$ for which $|J[\gamma]_x^M(q)-K(q)|\leq\beta_{|x|}(q)$, let $Q_>$ be the set of $q\in Q_x^M$ for which $J[\gamma_x]_x^M(q)>K(q)+\beta_{|x|}(q)$ and let $Q_<$ be the set of $q\in Q_x^M$ for which $J[\gamma_x](q)<K(q)-\beta_{|x|}(q)$. Define the oracle $K[x]$ as follows: $K[x](q)=J[\gamma_x](q)$ if $q\in Q_{\mathrm{close}}$ and $K[x](q)=K(q)$ otherwise.
By definition $K[x]$ is $\beta$-close to $K$ and therefore by hypothesis $M^{K[x]}$ computes $L$ with advantage $\epsilon$. In particular $M^{K[x]}$ on input $x$ differs from $L(x)$ with probability at most $(1-\epsilon(|x|))/2$. This completes the first part of the lemma.
For the second part of the lemma, assume that $x$ is such that $K(x)<(\beta(|x|)-C_1)/C_0$ (where $C_0,C_1$ are given in Lemma 17).
If all of the queries asked by $M$ on input $x$ are in $Q_{\mathrm{close}}$ then all query answers during the computation $M^{J[\gamma]}(x)$ are identical to those during the computation $M^{K[x]}(x)$, and therefore the output is the same.
So it now suffices to prove an upper bound of $\epsilon(x)/4$ on the probability that at least one of the selected queries on input $x$ belongs to $Q_<\cup Q_>$.
From the hypothesis, we have $K(x)\leq(\beta(|x|)-C_1)/C_0$. We now show that this implies that $Q_<=\varnothing$. By the first part of Lemma 17, for all queries $q\in Q_x^M$ we have $K(q)-J_x^M(q)\leq C_0K(x)+C_1$ which is at most $\beta(|x|)$ and so $q\notin Q_<$.
Hence we only need to upper bound the probability that, on input $x$, at least one query belongs to $Q_>$. For brevity let $k=k_x^M$. For $i\in\{1,\ldots,k\}$ let $p_i$ be the probability that the $i$th query is in $Q_>$. Let $\sigma=\sum_ip_i$, it suffices to show that $\sigma\leq\epsilon(|x|)/4$. Note that $\sigma/k$ is the chance that a random query $i$ (with respect to the random bits and the choice of $i\in\{1,\ldots,k\}$) is in $Q_>$. By the second part of Lemma 17, with $u=\beta(|x|)-2\log(K(q))$, $p\leq2^{2+2\log(K(q))-\beta(|x|)}$. Using the hypothesis on $\beta(|x|)$ this is at most $\epsilon(|x|)/4k$ and therefore $\sigma\leq\epsilon(|x|)/4$, as required.
\end{proof}
